\newpage
\section*{CONCLUSION}
\phantomsection\addcontentsline{toc}{section}{CONCLUSION}
\setcounter{section}{5}
\setcounter{subsection}{0}
\setcounter{subsubsection}{0}
\setcounter{paragraph}{0}
\setcounter{figure}{0}
\setcounter{table}{0}

%%%%%%%%%%%%%%%%%%%%%%%%%%%%%%%%%%%%%%%%%%%%%%%%%%%%%%%%%%%%%%%%%%%%%%%%%%%%%%%%%%%
\subsection*{Key Features}

The hardware--software co-design of a quantized HiFiC GAN accelerator for learned image compression has yielded positive results, demonstrating the capability of adversarial learning in optimizing both visual representation and reconstruction. A standout achievement of this research is the successful transition from algorithmic theory to a high-performance hardware implementation. Experimental results indicate that the hardware-accelerated generator processes one image in 735.6~ms, about $1459\times$ faster than a single-threaded CPU software reference (1073.19~s), while preserving the reconstruction quality of the original FP32 model after mixed-precision FP16 quantization.

The main features of the proposed design can be summarized as follows:
\begin{itemize}
    \item \textbf{Hardware--software co-design:} the complete HiFiC generator (the Fusion, UpConv, and Conv77 stages) is offloaded to a single custom FPGA IP, while the compression front-end and model orchestration remain in software on the Processing System.
    \item \textbf{Mixed-precision quantization:} the model is quantized to FP16 internally with FP32 interfaces, reducing memory footprint and computation cost with negligible loss in PSNR, SSIM, and LPIPS.
    \item \textbf{Efficient generator accelerator:} streaming line buffers, sliding-window generation, a parallel PE cluster with a rotating accumulator (sustaining $\text{II}=1$ in FP16), and shared convolution and normalization/activation engines minimize off-chip memory traffic and exploit the regular convolution structure of the generator.
    \item \textbf{Practical deployment:} the accelerator is implemented and evaluated on the ZCU104 platform at 300~MHz (timing closed at 308.74~MHz), processing one image in 735.6~ms at a throughput of 122.8~GOPS, about $1459\times$ faster than a single-threaded CPU software reference.
\end{itemize}

The project effectively bridges the gap between deep learning complexity and practical efficiency, showcasing a robust and adaptable solution for next-generation image processing. Rather than competing on raw wall-clock speed against high-end GPUs or peak-optimized CPU libraries, the accelerator is positioned for deterministic, low-latency inference within an embedded power envelope, where performance-per-watt rather than absolute latency is the decisive advantage.

%%%%%%%%%%%%%%%%%%%%%%%%%%%%%%%%%%%%%%%%%%%%%%%%%%%%%%%%%%%%%%%%%%%%%%%%%%%%%%%%%%%
\subsection*{Future Works}

Moving forward, the project will focus on several key trajectories to transition the proposed accelerator from a research prototype toward a robust, real-world solution:
\begin{itemize}
    \item \textbf{Hardware scaling:} extending the accelerator to higher input resolutions (beyond $256\times256$) and deepening inter-core pipelining so that the Fusion, UpConv, and Conv77 stages overlap, in order to raise end-to-end throughput.
    \item \textbf{Real-time deployment:} further optimizing the model for embedded and FPGA platforms, specifically targeting reduced computational complexity and lower latency.
    \item \textbf{Scalability:} extending the design to handle ultra-high-resolution datasets and diverse visual domains, ensuring robustness across different content types.
    \item \textbf{Hybrid modeling:} exploring the integration of GANs with Vision Transformers or Diffusion Models to push the boundaries of reconstruction fidelity.
    \item \textbf{Secure and efficient transmission:} integrating enhanced security measures and optimized bit-rate allocation, transforming the model into a versatile tool for high-efficiency, secure image transmission in bandwidth-constrained environments.
\end{itemize}
